\documentclass{aomart}

\usepackage[T1]{fontenc}
\usepackage[nopatch=footnote]{microtype}
\usepackage{amssymb,mathtools}
\usepackage{aliascnt}
\usepackage{enumitem}
\usepackage[nameinlink,noabbrev]{cleveref}

\allowdisplaybreaks
\hyphenpenalty=10000
\exhyphenpenalty=10000
\tolerance=2000
\emergencystretch=3em

\makeatletter
\def\@setsubjects{%
  \let\@subjects\@empty
  \ifx\@primarysubjects\@empty\else
    \edef\@subjects{Primary:~\@primarysubjects}%
  \fi
  \ifx\@secondarysubjects\@empty\else
    \ifx\@primarysubjects\@empty
      \edef\@subjects{Secondary:~\@secondarysubjects}%
    \else
      \edef\@subjects{\@subjects, Secondary:~\@secondarysubjects}%
    \fi
  \fi
  \ifpdf
    \hypersetup{pdfsubject=\@subjects}%
  \fi
  \if@aom@oldkeywords\else
    \@footnotetext{AMS Classification:~\@subjects.}%
  \fi
}
\makeatother

\setlist[enumerate]{topsep=4pt,itemsep=2pt}
\setlist[itemize]{topsep=4pt,itemsep=2pt}

\newtheorem{theorem}{Theorem}[section]
\newaliascnt{proposition}{theorem}
\newtheorem{proposition}[proposition]{Proposition}
\aliascntresetthe{proposition}
\newaliascnt{lemma}{theorem}
\newtheorem{lemma}[lemma]{Lemma}
\aliascntresetthe{lemma}
\newaliascnt{corollary}{theorem}
\newtheorem{corollary}[corollary]{Corollary}
\aliascntresetthe{corollary}
\theoremstyle{definition}
\newaliascnt{definition}{theorem}
\newtheorem{definition}[definition]{Definition}
\aliascntresetthe{definition}
\newaliascnt{example}{theorem}
\newtheorem{example}[example]{Example}
\aliascntresetthe{example}
\newaliascnt{standardfact}{theorem}
\newtheorem{standardfact}[standardfact]{Standard Fact}
\aliascntresetthe{standardfact}
\theoremstyle{remark}
\newaliascnt{remark}{theorem}
\newtheorem{remark}[remark]{Remark}
\aliascntresetthe{remark}

\crefname{theorem}{theorem}{theorems}
\crefname{proposition}{proposition}{propositions}
\crefname{lemma}{lemma}{lemmas}
\crefname{corollary}{corollary}{corollaries}
\crefname{definition}{definition}{definitions}
\crefname{example}{example}{examples}
\crefname{standardfact}{standard fact}{standard facts}
\crefname{remark}{remark}{remarks}

\newcommand{\bbk}{\Bbbk}
\newcommand{\OO}{\mathcal O}
\newcommand{\Acal}{\mathcal A}
\newcommand{\Bcal}{\mathcal B}
\newcommand{\Pcal}{\mathcal P}
\newcommand{\Qcal}{\mathcal Q}
\newcommand{\Scal}{\mathcal S}
\newcommand{\Rcal}{\mathcal R}
\newcommand{\End}{\operatorname{End}}
\newcommand{\im}{\operatorname{im}}
\newcommand{\Ker}{\operatorname{Ker}}
\newcommand{\coker}{\operatorname{coker}}
\newcommand{\Frac}{\operatorname{Frac}}
\newcommand{\length}{\operatorname{length}}
\newcommand{\val}{\operatorname{val}}
\newcommand{\wt}{\operatorname{wt}}
\newcommand{\Supp}{\operatorname{Supp}}
\newcommand{\id}{\mathrm{id}}
\newcommand{\diag}{\operatorname{diag}}
\newcommand{\Res}{\operatorname{Res}}
\newcommand{\lcmop}{\operatorname{lcm}}
\newcommand{\expn}{\operatorname{exp}}

\title[Refinement towers and integral separation]{Refinement Towers, Weighted Transfer, and Integral Separation of Operator Packets}
\author[Komorebi P.]{Komorebi P.}
\subject{primary}{msc2020}{47A68}
\subject{secondary}{msc2020}{13F30}
\subject{secondary}{msc2020}{15A21}
\subject{secondary}{msc2020}{05C20}
\subject{secondary}{msc2020}{42C40}
\keyword{idempotent refinement}
\keyword{filtered operator algebra}
\keyword{weighted paths}
\keyword{spectral projectors}
\keyword{Smith normal form}
\copyrightnote{}
\doinumber{}

\pdfstringdefDisableCommands{\let\unskip\relax}
\hypersetup{
  pdftitle={Refinement Towers, Weighted Transfer, and Integral Separation of Operator Packets},
  pdfauthor={Komorebi P.},
  pdfsubject={Nested idempotents, filtered inverses, lattice separation, and refinement defects}
}

\begin{document}
\begin{abstract}
A complete family of idempotents gives exact coordinates, but a single family does not record scale or integral visibility.  We study towers in which every coarse idempotent is the sum of its fine descendants.  The coordinate changes are invertible and coherent, and every operator block is the sum of its refined blocks.  Residue selectors give a model over arbitrary rings.  When roots of unity are available they agree with character averages, and coprime moduli factor through the Chinese remainder theorem.  In a complete filtered algebra, a diagonal operator plus positive weight transfers has an inverse obtained by summing admissible weighted paths.  The formula truncates exactly at finite precision, while finite transfer graphs have rational path series.  For a rational spectral decomposition over a discrete valuation ring, resultants give denominator bounds.  The exact denominator is the largest Smith exponent of the lattice gluing module.  This yields a necessary and sufficient precision delay, monotonicity under refinement, and a full congruence solving formula.  We close with exact formulas for the propagation of a coarse graining defect through powers, resolvents, holomorphic functions, exponentials, and multilevel towers.
\end{abstract}
\maketitle
\thispagestyle{plain}

\section{Introduction}

A complete family of pairwise annihilating idempotents decomposes a module into exact coordinate pieces.  Many problems require more than one such decomposition.  A signal may be split first into two residue classes and then into four.  A filtered operator may be resolved into transfers that appear at different valuation depths.  A rational spectral splitting may exist over the fraction field while remaining invisible on the original lattice.  These are all refinement problems, but they involve different notions of scale.

The first part of the article concerns exact refinement.  If a fine packet family refines a coarse family, then each coarse idempotent is the sum of its fine descendants.  The coordinate spaces at the two levels are canonically isomorphic because both represent the whole module.  The refinement relation adds more information.  It tells us which fine coordinates belong to each coarse coordinate, and it gives a compatible decomposition of every operator block.  This is the algebra behind polyphase coordinates and nested filter banks.  Standard treatments of those subjects may be found in \cite{Vaidyanathan1993,Daubechies1992}.

The second part studies inversion in a complete filtered algebra.  Suppose an operator has the form
\[
T=D+N,
\]
where $D$ is invertible and every packet transfer in $N$ raises the filtration by a positive amount.  The ordinary Neumann series can then be expanded into weighted paths.  The weight is not the number of edges.  It is the total filtration gain along the path.  At precision $L$, every path of weight at least $L$ vanishes.  A finite transfer graph also gives a rational generating function for the number of paths of each total weight.  The relevant generating function facts are classical and are discussed in \cite{Stanley2012}.

The third part concerns integral separation.  Let $V$ be a vector space over the fraction field of a discrete valuation ring $\OO$, and let $\Lambda$ be an $\OO$ lattice in $V$.  A rational projector $E$ need not preserve $\Lambda$.  Its separation weight is the least $s$ such that
\[
\pi^sE\Lambda\subseteq\Lambda.
\]
A Bézout construction gives an upper bound through a polynomial resultant.  That bound need not be exact because it does not see the position of the lattice.  The exact answer is controlled by the finite module
\[
G_{\Lambda,\Pcal}
=
\Lambda\Big/\bigoplus_{\alpha}(\Lambda\cap V_\alpha),
\]
where $V=\bigoplus_\alpha V_\alpha$ is the rational packet decomposition.  The largest Smith exponent of this gluing module is exactly the common denominator needed to make every packet projector integral.  This interpretation also proves that separation weight cannot decrease under refinement.  Background on local fields, Smith normal form, and matrix polynomials may be found in \cite{Serre1979,Newman1972,GohbergLancasterRodman2009}.

The last part records how a defect between fine and coarse operators propagates.  If
\[
\Delta=CT_f-T_cC,
\]
then exact formulas describe the corresponding defect for powers, resolvents, holomorphic functions, exponentials, and a chain of several refinement levels.  These formulas are elementary, but collecting them in one place makes the later analytic uses transparent.  The spectral and resolvent background follows the usual operator theory in \cite{Kato1995}.

The paper remains linear.  Exact refinement is not approximation, and no packet is discarded.  Nonlinear operation trees, nonlinear correction, and obstruction complexes require a separate theory.

\section{Four kinds of depth}\label{sec:four_depths}

Several quantities below are measured by nonnegative integers.  They do not measure the same thing.  Keeping them separate makes the later formulas easier to read.

The first is packet depth.  A finer packet family has more idempotents and a more detailed address set.  The passage from residues modulo two to residues modulo four is an example.  Packet depth is combinatorial.  It says how many distinctions have been made in the coordinate system.

The second is filtration depth.  A transfer in $F^r\Bcal$ is invisible modulo $F^r\Bcal$ and becomes visible at greater precision.  Two paths with the same number of edges can have different filtration depths.  A longer path can appear earlier when its edges have smaller weights.

The third is separation depth.  A rational packet projector may require multiplication by $\pi^s$ before it preserves the chosen lattice.  This number depends on both the rational decomposition and the lattice.  It can be zero even when the corresponding eigenvalues are close, and it can be positive even when the rational projector is simple.

The fourth is compatibility depth.  A fine operator and a separately chosen coarse operator need not commute with coarsening.  Their discrepancy is
\[
\Delta=CT_f-T_cC.
\]
This is not a packet depth or a denominator.  It measures the failure of one operator model to represent another exactly.

The four notions meet in applications, but no one of them determines the others.  An exact refinement tower may have zero coordinate error while its fine spectral packets have positive lattice weight.  A filtered inverse may truncate exactly at a given depth while the coarse operator still has a nonzero compatibility defect.  The paper treats each quantity with its own theorem before combining them.

A simple two packet model already contains all four ideas.  Let
\[
X=X_0\oplus X_1,
\qquad
D=\begin{pmatrix}D_0&0\\0&D_1\end{pmatrix},
\qquad
N=\begin{pmatrix}0&N_{01}\\N_{10}&0\end{pmatrix}.
\]
A finer family may split $X_0$ and $X_1$ into smaller summands.  The blocks $N_{01}$ and $N_{10}$ may enter at different filtration weights.  Over a fraction field, the rational packet spaces may fail to split the chosen lattice.  A coarse operator on one selected summand may also fail to intertwine with the full matrix.  The later sections isolate these four effects one at a time.

\section{Packet systems and refinement towers}\label{sec:refinement}

We begin with finite packet systems.  No inner product is needed.

\begin{definition}[Packet system]\label{def:packet_system}
Let $X$ be a module over a unital ring.  A finite packet system on $X$ is a family
\[
\Pcal=(P_\alpha)_{\alpha\in I}
\subseteq\End(X)
\]
of nonzero idempotents such that
\[
P_\alpha P_\beta=0
\quad\text{when }\alpha\ne\beta,
\qquad
\sum_{\alpha\in I}P_\alpha=I_X.
\]
Write
\[
X_\alpha=P_\alpha X,
\qquad
X^{\Pcal}=\bigoplus_{\alpha\in I}X_\alpha.
\]
The analysis and synthesis maps are
\[
\Acal_{\Pcal}x=(P_\alpha x)_\alpha,
\qquad
\Scal_{\Pcal}(x_\alpha)_\alpha=\sum_\alpha x_\alpha.
\]
They are inverse isomorphisms.
\end{definition}

\begin{definition}[Refinement]\label{def:refinement}
Let $\Pcal=(P_\alpha)_{\alpha\in I}$ and $\Qcal=(Q_\beta)_{\beta\in J}$ be packet systems on the same module.  We say that $\Qcal$ refines $\Pcal$ when there is a surjection
\[
\rho:J\longrightarrow I
\]
such that
\[
P_\alpha=\sum_{\rho(\beta)=\alpha}Q_\beta
\]
for every $\alpha\in I$.  The index $\beta$ is then called a child of $\rho(\beta)$.
\end{definition}

\begin{theorem}[Children, parents, and transitivity]\label{thm:children_transitivity}
Suppose that $\Qcal$ refines $\Pcal$ through $\rho$.  Then
\[
P_\alpha Q_\beta
=
Q_\beta P_\alpha
=
\begin{cases}
Q_\beta & \rho(\beta)=\alpha,\\
0 & \rho(\beta)\ne\alpha.
\end{cases}
\]
Hence $Q_\beta X\subseteq P_{\rho(\beta)}X$.  If a third packet system $\Rcal$ refines $\Qcal$, then $\Rcal$ refines $\Pcal$ through the composite parent map.
\end{theorem}

\begin{proof}
Multiply the identity
\[
P_\alpha=\sum_{\rho(\gamma)=\alpha}Q_\gamma
\]
by $Q_\beta$ on either side.  Pairwise annihilation in the fine family leaves $Q_\beta$ when $\rho(\beta)=\alpha$ and leaves zero otherwise.  For transitivity, substitute the child expansion of each $Q_\beta$ into the expansion of $P_\alpha$ and group the resulting terms by the composite parent.
\end{proof}

A refinement tower is a sequence $(\Pcal_\ell)_{\ell\geq0}$ with parent maps
\[
\rho_{m,\ell}:I_m\longrightarrow I_\ell
\qquad(m\geq\ell)
\]
that compose as expected.  The next result shows that the coordinate changes form a coherent system.

\begin{definition}[Coordinate change across levels]\label{def:coordinate_change}
For two packet systems $\Pcal$ and $\Qcal$ on $X$, define
\[
J_{\Qcal\leftarrow\Pcal}
=
\Acal_{\Qcal}\Scal_{\Pcal}:X^{\Pcal}\longrightarrow X^{\Qcal}
\]
and
\[
C_{\Pcal\leftarrow\Qcal}
=
\Acal_{\Pcal}\Scal_{\Qcal}:X^{\Qcal}\longrightarrow X^{\Pcal}.
\]
When $\Qcal$ refines $\Pcal$, the first map resolves each coarse coordinate into its children and the second adds the children back to their parent.
\end{definition}

\begin{theorem}[Exact coordinate change and coherence]\label{thm:coordinate_coherence}
For any two packet systems on $X$,
\[
C_{\Pcal\leftarrow\Qcal}J_{\Qcal\leftarrow\Pcal}=I_{X^{\Pcal}},
\qquad
J_{\Qcal\leftarrow\Pcal}C_{\Pcal\leftarrow\Qcal}=I_{X^{\Qcal}}.
\]
For three packet systems $\Pcal$, $\Qcal$, and $\Rcal$ on $X$,
\[
J_{\Rcal\leftarrow\Qcal}J_{\Qcal\leftarrow\Pcal}
=
J_{\Rcal\leftarrow\Pcal},
\]
\[
C_{\Pcal\leftarrow\Qcal}C_{\Qcal\leftarrow\Rcal}
=
C_{\Pcal\leftarrow\Rcal}.
\]
If $\Qcal$ refines $\Pcal$ through $\rho$, then
\[
\bigl(J_{\Qcal\leftarrow\Pcal}x\bigr)_\beta
=
Q_\beta x_{\rho(\beta)},
\qquad
\bigl(C_{\Pcal\leftarrow\Qcal}y\bigr)_\alpha
=
\sum_{\rho(\beta)=\alpha}y_\beta.
\]
\end{theorem}

\begin{proof}
The packet analysis and synthesis maps are inverse.  Therefore
\[
C_{\Pcal\leftarrow\Qcal}J_{\Qcal\leftarrow\Pcal}
=
\Acal_{\Pcal}\Scal_{\Qcal}\Acal_{\Qcal}\Scal_{\Pcal}
=
\Acal_{\Pcal}\Scal_{\Pcal}
=
I.
\]
The other inverse identity is the same calculation with the two systems exchanged.  The coherence identities follow by inserting
\[
\Scal_{\Qcal}\Acal_{\Qcal}=I_X.
\]
Now assume that $\Qcal$ refines $\Pcal$.  If $x=(x_\alpha)_\alpha$ with $x_\alpha\in P_\alpha X$, then
\[
Q_\beta\sum_\alpha x_\alpha
=
Q_\beta x_{\rho(\beta)}
\]
by \cref{thm:children_transitivity}.  The formula for $C$ follows by synthesis followed by projection onto the parent packet.
\end{proof}

\begin{remark}[Refinement and reduction]\label{rem:refinement_reduction}
The maps in \cref{thm:coordinate_coherence} are isomorphisms because both coordinate spaces represent all of $X$.  They do not reduce dimension and they do not discard packets.  Approximation begins only after one chooses a truncation or compression rule.
\end{remark}

\begin{corollary}[Orthogonal Hilbert packets]\label{cor:orthogonal_isometry}
If $X$ is a Hilbert space and every packet system in the tower consists of orthogonal projections, then analysis and synthesis are unitary.  Every coordinate change in \cref{thm:coordinate_coherence} is unitary.
\end{corollary}

\begin{proof}
For an orthogonal packet system,
\[
\|x\|^2=\sum_\alpha\|P_\alpha x\|^2.
\]
Thus analysis is an isometry onto the full packet direct sum.  Synthesis is its inverse.  A coordinate change is a composition of two unitary maps.
\end{proof}

\subsection{Operator blocks across refinement}

For $T\in\End(X)$ and a packet system $\Pcal$, write
\[
T^{\Pcal}_{\beta\alpha}=P_\beta T P_\alpha.
\]

\begin{theorem}[Refinement of operator blocks]\label{thm:block_refinement}
Suppose that $\Qcal$ refines $\Pcal$ through $\rho$.  Then
\[
T^{\Pcal}_{ba}
=
\sum_{\substack{\rho(\beta)=b\\\rho(\alpha)=a}}
T^{\Qcal}_{\beta\alpha}.
\]
The matrix of $T$ in fine coordinates is related to the matrix in coarse coordinates by
\[
[T]_{\Pcal}
=
C_{\Pcal\leftarrow\Qcal}[T]_{\Qcal}J_{\Qcal\leftarrow\Pcal}.
\]
\end{theorem}

\begin{proof}
Expand the two coarse projectors into their children:
\[
P_bTP_a
=
\left(\sum_{\rho(\beta)=b}Q_\beta\right)
T
\left(\sum_{\rho(\alpha)=a}Q_\alpha\right).
\]
Distributing the finite sums gives the first identity.  For the coordinate identity, use
\[
[T]_{\Pcal}=\Acal_{\Pcal}T\Scal_{\Pcal}
\]
and insert
\[
I_X=\Scal_{\Qcal}\Acal_{\Qcal}.
\]
\end{proof}

\begin{corollary}[Support under refinement]\label{cor:block_support_refinement}
Under the hypotheses of \cref{thm:block_refinement}, a coarse block $T^{\Pcal}_{ba}$ is zero if and only if every fine descendant block
\[
T^{\Qcal}_{\beta\alpha}
\qquad
\bigl(\rho(\beta)=b,\ \rho(\alpha)=a\bigr)
\]
is zero.
\end{corollary}

\begin{proof}
If every fine block is zero, the coarse block is zero by \cref{thm:block_refinement}.  Conversely, suppose $P_bTP_a=0$.  Multiply on the left by $Q_\beta$ and on the right by $Q_\alpha$.  The child identities in \cref{thm:children_transitivity} give
\[
Q_\beta TQ_\alpha=0.
\]
\end{proof}

Thus refinement cannot hide an active transfer inside a zero coarse block.  It can only reveal where the active transfer sits among the descendants.

\subsection{Common refinements}

Two packet systems need not be comparable.  When their idempotents commute, their products give a common refinement.

\begin{theorem}[Common refinement of commuting packet systems]\label{thm:common_refinement}
Let $\Pcal=(P_\alpha)_{\alpha\in I}$ and $\Qcal=(Q_\beta)_{\beta\in J}$ be packet systems on $X$.  Assume
\[
P_\alpha Q_\beta=Q_\beta P_\alpha
\]
for all $\alpha$ and $\beta$.  Remove the zero operators from the family
\[
R_{\alpha\beta}=P_\alpha Q_\beta.
\]
The remaining operators form a packet system that refines both $\Pcal$ and $\Qcal$.  The same statement holds for any finite collection of pairwise commuting packet systems.
\end{theorem}

\begin{proof}
Commutativity gives
\[
R_{\alpha\beta}^2=P_\alpha^2Q_\beta^2=R_{\alpha\beta}.
\]
If $(\alpha,\beta)\ne(\alpha',\beta')$, then either $P_\alpha P_{\alpha'}=0$ or $Q_\beta Q_{\beta'}=0$, so the corresponding products annihilate one another.  Finally,
\[
\sum_{\alpha,\beta}R_{\alpha\beta}
=
\left(\sum_\alpha P_\alpha\right)
\left(\sum_\beta Q_\beta\right)
=I_X.
\]
For a fixed $\alpha$,
\[
P_\alpha
=P_\alpha\sum_\beta Q_\beta
=
\sum_\beta R_{\alpha\beta},
\]
and the same argument works for a fixed $\beta$.  Iteration proves the finite family statement.
\end{proof}

\begin{proposition}[Coarsest common refinement]\label{prop:coarsest_common_refinement}
Under the hypotheses of \cref{thm:common_refinement}, let $\Rcal$ be a packet system that refines both $\Pcal$ and $\Qcal$.  Then $\Rcal$ refines the nonzero product family $(P_\alpha Q_\beta)$.  Hence that product family is the coarsest packet system that refines both $\Pcal$ and $\Qcal$.
\end{proposition}

\begin{proof}
Let $R_\gamma$ be a packet in $\Rcal$.  Because $\Rcal$ refines $\Pcal$, there is a unique $\alpha$ with
\[
P_\alpha R_\gamma=R_\gamma.
\]
Likewise there is a unique $\beta$ with
\[
Q_\beta R_\gamma=R_\gamma.
\]
Therefore
\[
P_\alpha Q_\beta R_\gamma=R_\gamma.
\]
Every packet of $\Rcal$ lies in one product packet.  Summing the packets with the same pair of parents reconstructs that product packet.
\end{proof}

\begin{example}[A two level split]\label{ex:two_level_split}
Let
\[
X=X_{00}\oplus X_{01}\oplus X_{10}\oplus X_{11}.
\]
The four coordinate projections form a fine packet system.  The two sums
\[
P_0=Q_{00}+Q_{01},
\qquad
P_1=Q_{10}+Q_{11}
\]
form a coarse packet system.  If $x=(x_{00},x_{01},x_{10},x_{11})$, the coarse coordinates are
\[
(x_{00}+x_{01},x_{10}+x_{11})
\]
when the sums are interpreted inside the two coarse direct summands.  Fine analysis recovers the four entries exactly.  Nothing has been approximated.
\end{example}

\section{Residue coordinates and coprime factorization}\label{sec:residue}

Residue packets provide a model that works over every coefficient ring.  The character formula is a second description that needs stronger hypotheses.

\begin{definition}[Residue selector]\label{def:residue_selector}
Let $R$ be a commutative ring and let
\[
F(z)=\sum_{n\geq0}a_nz^n\in R[[z]].
\]
For $N\geq1$ and $0\leq j<N$, define
\[
(P_{N,j}F)(z)
=
\sum_{n\equiv j\pmod N}a_nz^n.
\]
\end{definition}

\begin{proposition}[Residue packet algebra]\label{prop:residue_algebra}
The operators $(P_{N,j})_{0\leq j<N}$ form a packet system on $R[[z]]$:
\[
P_{N,i}P_{N,j}=\delta_{ij}P_{N,j},
\qquad
\sum_{j=0}^{N-1}P_{N,j}=I.
\]
\end{proposition}

\begin{proof}
Every nonnegative integer belongs to one residue class modulo $N$.  A coefficient survives both $P_{N,i}$ and $P_{N,j}$ exactly when the two labels are equal to its residue.  Summing all selectors preserves every coefficient once.
\end{proof}

\begin{theorem}[Character realization]\label{thm:character_realization}
Let $\bbk$ be a field, assume that $N$ is nonzero in $\bbk$, and suppose that $\bbk$ contains a primitive $N$th root of unity $\omega$.  Define
\[
(UF)(z)=F(\omega z).
\]
Then
\[
P_{N,j}
=
\frac1N\sum_{k=0}^{N-1}\omega^{-jk}U^k.
\]
\end{theorem}

\begin{proof}
On $z^n$, the right side acts by the scalar
\[
\frac1N\sum_{k=0}^{N-1}\omega^{k(n-j)}.
\]
This is one when $n\equiv j\pmod N$ and zero otherwise by the geometric sum identity in the field $\bbk$.
\end{proof}

\begin{remark}[Two distinct constructions]\label{rem:selector_character}
The coefficient selector of \cref{def:residue_selector} uses no division.  The character formula in \cref{thm:character_realization} divides by $N$.  At a prime that divides $N$, the coefficient selector may remain integral even when the character average does not.
\end{remark}

\begin{theorem}[Nested residue refinement]\label{thm:residue_refinement}
If $N$ divides $M$, then
\[
P_{N,j}
=
\sum_{\substack{0\leq s<M\\s\equiv j\pmod N}}P_{M,s}.
\]
Thus the packet system modulo $M$ refines the packet system modulo $N$.
\end{theorem}

\begin{proof}
An exponent $n$ satisfies $n\equiv j\pmod N$ exactly when its unique residue $s$ modulo $M$ satisfies $s\equiv j\pmod N$.
\end{proof}

\begin{example}[The shift graph of residue packets]\label{ex:residue_shift}
Multiplication by $z$ sends the residue packet $P_{N,j}R[[z]]$ into $P_{N,j+1}R[[z]]$, where the index is read modulo $N$.  Its packet graph is the directed cycle
\[
0\longrightarrow1\longrightarrow\cdots\longrightarrow N-1\longrightarrow0.
\]
Multiplication by a polynomial supported on a set of residues gives the corresponding finite set of phase shifts.  This is the simplest transfer graph behind a polyphase representation.
\end{example}

\begin{theorem}[Coprime factorization]\label{thm:crt_factorization}
Let $A$ and $B$ be coprime positive integers.  For each $j$ modulo $AB$, write
\[
j_A=j\pmod A,
\qquad
j_B=j\pmod B.
\]
Then
\[
P_{AB,j}=P_{A,j_A}P_{B,j_B}.
\]
Every pair $(j_A,j_B)$ occurs for one residue $j$ modulo $AB$.
\end{theorem}

\begin{proof}
A coefficient survives the product on the right exactly when its exponent satisfies both congruences.  The Chinese remainder theorem gives one residue class modulo $AB$ with those two coordinates.  The needed form of the theorem is standard and may be found in \cite{Lang2002}.
\end{proof}

\begin{corollary}[Dyadic and odd coordinates]\label{cor:dyadic_odd_coordinates}
Write $N=2^aM$ with $M$ odd.  Every residue packet modulo $N$ is the product of one packet modulo $2^a$ and one packet modulo $M$.
\end{corollary}

\begin{proof}
Apply \cref{thm:crt_factorization} to the coprime moduli $2^a$ and $M$.
\end{proof}

\begin{example}[Residues modulo twelve]\label{ex:mod_twelve}
The pair
\[
(j\pmod4,j\pmod3)=(3,1)
\]
corresponds to $j=7$ modulo $12$.  Hence
\[
P_{12,7}=P_{4,3}P_{3,1}.
\]
This factorization concerns addresses.  It does not say that the two factors have the same integral spectral behavior.
\end{example}

\section{Complete filtered transfer algebras}\label{sec:filtered}

Weighted inversion is most transparent in an abstract multiplicative filtration.  A discrete valuation ring is the main example.

\begin{definition}[Complete multiplicative filtration]\label{def:filtered_algebra}
A complete filtered algebra is a unital algebra $\Bcal$ with a decreasing sequence
\[
\Bcal=F^0\Bcal\supseteq F^1\Bcal\supseteq F^2\Bcal\supseteq\cdots
\]
such that
\[
F^r\Bcal\,F^s\Bcal\subseteq F^{r+s}\Bcal,
\qquad
\bigcap_{L\geq0}F^L\Bcal=0,
\]
and the natural map
\[
\Bcal\longrightarrow\varprojlim_L\Bcal/F^L\Bcal
\]
is an isomorphism.  A transfer term $N_e$ has declared weight $w(e)\geq1$ when
\[
N_e\in F^{w(e)}\Bcal.
\]
\end{definition}

If $\OO$ is a complete discrete valuation ring with uniformizer $\pi$, the filtration
\[
F^L\Bcal=\pi^L\Bcal
\]
is the standard example.  Basic facts about completeness in this setting are reviewed in \cite{Serre1979}.

\begin{theorem}[Weighted inverse]\label{thm:weighted_inverse}
Let $\Bcal$ be a complete filtered algebra.  Let $D\in\Bcal$ be invertible with
\[
D,D^{-1}\in F^0\Bcal.
\]
Let $E$ be a finite set and suppose
\[
T=D+\sum_{e\in E}N_e,
\qquad
N_e\in F^{w(e)}\Bcal,
\qquad
w(e)\geq1.
\]
Set
\[
A_e=D^{-1}N_e.
\]
For a word $\gamma=(e_1,\ldots,e_r)$, write
\[
A_\gamma=A_{e_r}\cdots A_{e_1},
\qquad
W(\gamma)=\sum_{j=1}^rw(e_j).
\]
The empty word has weight zero and contributes the identity.  Then $T$ is invertible and
\[
T^{-1}
=
\sum_\gamma(-1)^{|\gamma|}A_\gamma D^{-1}.
\]
The series converges in the filtration, and the term indexed by $\gamma$ belongs to $F^{W(\gamma)}\Bcal$.
\end{theorem}

\begin{proof}
Put
\[
A=\sum_{e\in E}A_e.
\]
Since left multiplication by $D^{-1}$ preserves the filtration, $A\in F^1\Bcal$.  Completeness gives the convergent Neumann series
\[
(I+A)^{-1}=\sum_{r\geq0}(-A)^r.
\]
Because $T=D(I+A)$,
\[
T^{-1}=(I+A)^{-1}D^{-1}.
\]
Expanding each power of $A$ gives the sum over words, with the rightmost edge acting first.  Multiplicativity of the filtration gives
\[
A_\gamma\in F^{W(\gamma)}\Bcal.
\]
For each $L$, only finitely many words have weight below $L$ because $E$ is finite and every edge has positive weight.  Thus the word expansion converges to the same inverse.
\end{proof}

The preceding formula does not yet use packet blocks.  When $D$ is packet diagonal and the terms $N_e$ are individual transfers, the words become directed paths.

\begin{theorem}[Block path formula]\label{thm:block_path_formula}
Let $\Pcal=(P_\alpha)_{\alpha\in I}$ be a finite packet system in a complete filtered algebra of operators.  Suppose
\[
D=\sum_\alpha P_\alpha D P_\alpha
\]
is invertible on every packet, and write
\[
N=\sum_{e:\alpha\to\beta}N_e,
\qquad
N_e=P_\beta N_eP_\alpha,
\qquad
N_e\in F^{w(e)}\Bcal.
\]
For an edge $e:\alpha\to\beta$, put
\[
A_e=(P_\beta DP_\beta)^{-1}N_e.
\]
Then the block from $P_\alpha X$ to $P_\beta X$ of $(D+N)^{-1}$ is
\[
P_\beta(D+N)^{-1}P_\alpha
=
\sum_{\gamma:\alpha\leadsto\beta}
(-1)^{|\gamma|}A_\gamma D_\alpha^{-1},
\]
where $D_\alpha=P_\alpha DP_\alpha$ and the sum runs over directed paths from $\alpha$ to $\beta$, including the empty path when $\alpha=\beta$.
\end{theorem}

\begin{proof}
The diagonal inverse is
\[
D^{-1}=\sum_\alpha D_\alpha^{-1}P_\alpha.
\]
The block $P_\beta D^{-1}NP_\alpha$ is nonzero only when there is a declared edge from $\alpha$ to $\beta$.  Repeated block multiplication shows that the $r$th power of $D^{-1}N$ is the sum over all directed paths of length $r$, with the operator factors ordered along the path.  Apply \cref{thm:weighted_inverse} and multiply on the right by the packet diagonal inverse.
\end{proof}

\begin{corollary}[Minimum path weight]\label{cor:minimum_path_weight}
For packets $\alpha$ and $\beta$, let $d(\alpha,\beta)$ be the least total weight of a directed path from $\alpha$ to $\beta$, with $d=\infty$ when no path exists.  If $\alpha\ne\beta$, then
\[
P_\beta(D+N)^{-1}P_\alpha
\in
F^{d(\alpha,\beta)}\Bcal.
\]
If there is no path from $\alpha$ to $\beta$, the block is zero.
\end{corollary}

\begin{proof}
Every term in the block path formula is indexed by a path from $\alpha$ to $\beta$ and lies in the filtration level given by its total weight.  If no such path exists, the sum is empty.
\end{proof}

\begin{example}[A longer path that appears earlier]\label{ex:weighted_route}
Consider three packets $\alpha$, $\beta$, and $\gamma$.  Suppose there is a direct edge from $\alpha$ to $\gamma$ of weight three, together with edges
\[
\alpha\longrightarrow\beta,
\qquad
\beta\longrightarrow\gamma
\]
of weight one each.  The two edge route has total weight two, so it contributes to the inverse before the direct edge does.  This is why the truncation rule uses total weight rather than path length.  The lower bound in \cref{cor:minimum_path_weight} need not be an equality because several paths of the same weight may cancel.
\end{example}

\begin{corollary}[Acyclic transfer graph]\label{cor:acyclic_inverse}
If the packet transfer graph is acyclic, then $D^{-1}N$ is nilpotent.  The inverse in \cref{thm:block_path_formula} is a finite sum and does not require completeness of the filtration.
\end{corollary}

\begin{proof}
A finite acyclic directed graph has a bound on path length.  Every sufficiently long product of transfer blocks is therefore zero.  The Neumann series terminates.
\end{proof}

\begin{corollary}[Stability of inversion at finite precision]\label{cor:inverse_stability}
Let $T$ and $\widetilde T$ be invertible elements of $F^0\Bcal$.  If
\[
T-\widetilde T\in F^L\Bcal,
\]
then
\[
T^{-1}-\widetilde T^{-1}\in F^L\Bcal.
\]
\end{corollary}

\begin{proof}
Use the resolvent identity at zero in the form
\[
T^{-1}-\widetilde T^{-1}
=
T^{-1}(\widetilde T-T)\widetilde T^{-1}.
\]
Both inverse factors belong to $F^0\Bcal$, so the product lies in $F^L\Bcal$.
\end{proof}

\begin{corollary}[Exact finite precision truncation]\label{cor:finite_precision}
Under the hypotheses of \cref{thm:weighted_inverse}, define
\[
R_L
=
\sum_{W(\gamma)<L}(-1)^{|\gamma|}A_\gamma D^{-1}.
\]
Then
\[
T^{-1}-R_L\in F^L\Bcal,
\]
\[
TR_L-I\in F^L\Bcal,
\qquad
R_LT-I\in F^L\Bcal.
\]
Thus $R_L$ is the exact inverse in $\Bcal/F^L\Bcal$.
\end{corollary}

\begin{proof}
Every omitted word has weight at least $L$, so the difference from the full inverse lies in $F^L\Bcal$.  Multiplication by $T$ or by its inverse preserves membership in $F^L\Bcal$ under the stated filtration assumptions.  The two residual statements follow.
\end{proof}

\begin{example}[A triangular transfer]\label{ex:triangular_transfer}
Let $\OO$ be complete with uniformizer $\pi$, and consider
\[
T=
\begin{pmatrix}
1&\pi^a\\
0&1
\end{pmatrix}
\]
with $a\geq1$.  The only nontrivial transfer has weight $a$, and its square is zero.  Hence
\[
T^{-1}
=
\begin{pmatrix}
1&-\pi^a\\
0&1
\end{pmatrix}.
\]
At precision $L\leq a$, the transfer vanishes and the inverse is the identity modulo $\pi^L$.  At depth $a+1$, the off diagonal packet first becomes visible.
\end{example}

\section{Rational weighted transfer series}\label{sec:path_series}

A finite weighted graph records which path terms can occur in the inverse.  Its path counts have a rational generating function.

\begin{definition}[Weighted transfer series]\label{def:transfer_series}
Let $G$ be a finite directed multigraph.  Give every edge $e$ a positive integer weight $w(e)$ and a scalar multiplicity $m(e)$ in a field $\bbk$.  Define the polynomial adjacency matrix
\[
A(t)_{ij}
=
\sum_{e:i\to j}m(e)t^{w(e)}.
\]
For column vectors $u$ and $v$, define
\[
Z_{u,v}(t)
=
\sum_{n\geq0}a_nt^n,
\]
where $a_n$ is the total multiplicity of paths from the vertices selected by $u$ to the vertices selected by $v$ with total weight $n$.
\end{definition}

\begin{theorem}[Rationality of the transfer series]\label{thm:transfer_rationality}
The weighted transfer series satisfies
\[
Z_{u,v}(t)
=
u^{\mathsf T}(I-A(t))^{-1}v.
\]
It is a rational function whose denominator divides
\[
\det(I-A(t)).
\]
\end{theorem}

\begin{proof}
Since every entry of $A(t)$ has zero constant term, the inverse exists in the formal power series ring and
\[
(I-A(t))^{-1}=\sum_{r\geq0}A(t)^r.
\]
The $(i,j)$ entry of $A(t)^r$ is the sum of the monomials attached to all paths of length $r$ from $i$ to $j$.  Summing over $r$ records paths of every length.  Multiplication by the endpoint vectors gives the stated series.  The adjugate formula gives
\[
(I-A(t))^{-1}
=
\frac{\operatorname{adj}(I-A(t))}{\det(I-A(t))},
\]
so every entry is rational.
\end{proof}

\begin{corollary}[Recurrence for path counts]\label{cor:path_recurrence}
Write
\[
Z_{u,v}(t)=\frac{P(t)}{Q(t)},
\qquad
Q(t)=1+q_1t+\cdots+q_dt^d.
\]
Then, for every $n>\deg P$,
\[
a_n+q_1a_{n-1}+\cdots+q_da_{n-d}=0.
\]
\end{corollary}

\begin{proof}
Compare coefficients in the identity
\[
Q(t)Z_{u,v}(t)=P(t).
\]
Beyond the degree of $P$, the coefficient on the right is zero.
\end{proof}

\begin{corollary}[Several independent weights]\label{cor:multigraded_series}
Suppose every edge has a weight vector
\[
w(e)=(w_1(e),\ldots,w_r(e))\in\mathbb N^r\setminus\{0\}.
\]
Replace $t^{w(e)}$ by
\[
t_1^{w_1(e)}\cdots t_r^{w_r(e)}.
\]
The resulting multigraded transfer series is rational in $t_1,\ldots,t_r$.
\end{corollary}

\begin{proof}
The same matrix inverse proof works over the formal power series ring in several variables.  The adjugate formula gives rationality.
\end{proof}

The multigraded form is useful when dyadic depth, odd depth, and another filtration are tracked separately rather than collapsed into one number.

\begin{example}[A transfer series with one loop]\label{ex:simple_transfer_series}
Take vertices $0$, $1$, and $2$.  Give the edge $0\to1$ weight one, the loop $1\to1$ weight two, and the edge $1\to2$ weight one.  Every path from $0$ to $2$ uses the first and last edges and may traverse the loop any number of times.  Hence
\[
Z_{0,2}(t)
=
t^2(1+t^2+t^4+\cdots)
=
\frac{t^2}{1-t^2}.
\]
There is one path at every positive even weight and no path at an odd weight.  The denominator records the increase of two units produced by the loop.
\end{example}

\section{Rational spectral projectors and resultant bounds}\label{sec:resultant}

Let $\OO$ be a discrete valuation ring with fraction field $K$ and uniformizer $\pi$.  Let $\Lambda$ be a finite free $\OO$ module and put
\[
V=K\otimes_{\OO}\Lambda.
\]
Suppose $A\in\End_{\OO}(\Lambda)$.

\begin{theorem}[Bézout spectral projector]\label{thm:bezout_projector}
Assume that
\[
\chi_A(x)=f(x)g(x)
\]
with monic polynomials $f,g\in\OO[x]$ that are coprime in $K[x]$.  Choose $u,v\in\OO[x]$ and a nonzero $r\in\OO$ such that
\[
u(x)f(x)+v(x)g(x)=r.
\]
Then
\[
E_f=r^{-1}v(A)g(A)
\]
is the projector of $V$ onto $\Ker f(A)$ along $\Ker g(A)$.  In particular,
\[
E_f^2=E_f,
\qquad
rE_f\in\End_{\OO}(\Lambda).
\]
\end{theorem}

\begin{proof}
Cayley Hamilton gives
\[
f(A)g(A)=0.
\]
If $x\in\Ker f(A)$, then the Bézout identity gives
\[
v(A)g(A)x=rx,
\]
so $E_fx=x$.  If $x\in\Ker g(A)$, then $E_fx=0$.  The same identity gives the direct sum decomposition
\[
V=\Ker f(A)\oplus\Ker g(A).
\]
Indeed, every $x\in V$ is the sum of
\[
r^{-1}v(A)g(A)x\in\Ker f(A)
\]
and
\[
r^{-1}u(A)f(A)x\in\Ker g(A).
\]
Thus $E_f$ is the stated projector and is idempotent.  The formula for $rE_f$ has coefficients in $\OO$, so it preserves $\Lambda$.
\end{proof}

\begin{theorem}[Several coprime spectral packets]\label{thm:several_projectors}
Assume
\[
\chi_A(x)=\prod_{i=1}^rf_i(x),
\]
where the monic factors $f_i\in\OO[x]$ are pairwise coprime in $K[x]$.  Put
\[
g_i=\prod_{j\ne i}f_j.
\]
There are complementary projectors $E_i$ with
\[
\im E_i=\Ker f_i(A),
\qquad
\sum_iE_i=I,
\qquad
E_iE_j=0
\quad(i\ne j).
\]
One may choose $E_i$ so that
\[
\Res(f_i,g_i)E_i\in\End_{\OO}(\Lambda)
\]
up to multiplication of the resultant by a unit.
\end{theorem}

\begin{proof}
Apply \cref{thm:bezout_projector} to the factorization $\chi_A=f_ig_i$.  The resulting projector is one on $\Ker f_i(A)$ and zero on every $\Ker f_j(A)$ with $j\ne i$.  The direct sum decomposition follows from the polynomial Chinese remainder theorem.  The orthogonality and sum identities follow by evaluating the corresponding idempotents in $K[x]/(\chi_A)$.  The Sylvester Bézout relation for $(f_i,g_i)$ gives the stated denominator bound.
\end{proof}

\begin{corollary}[Resultant denominator bound]\label{cor:resultant_bound}
Under the hypotheses of \cref{thm:bezout_projector}, one may take
\[
r=\Res(f,g)
\]
up to multiplication by a unit.  Therefore
\[
\pi^{\val_\pi(\Res(f,g))}E_f\Lambda\subseteq\Lambda.
\]
\end{corollary}

\begin{proof}
The Sylvester matrix construction gives a Bézout relation with right side equal to the resultant.  Apply \cref{thm:bezout_projector}.  Standard facts about resultants may be found in \cite{Lang2002,GohbergLancasterRodman2009}.
\end{proof}

\begin{example}[A resultant bound that is not exact]\label{ex:loose_resultant}
Let
\[
A=\diag(0,\pi^s)
\]
on the standard lattice $\Lambda=\OO^2$.  The factors $x$ and $x-\pi^s$ have resultant $\pi^s$, but the two spectral projectors are the integral coordinate projectors.  Their exact separation weight is zero.  The resultant sees the closeness of the two eigenvalues.  It does not see that the chosen lattice already splits along their eigenspaces.
\end{example}

\section{Exact lattice separation}\label{sec:lattice}

The exact denominator belongs to the pair consisting of the rational packet decomposition and the integral lattice.

\begin{definition}[Lattice packet data]\label{def:lattice_packet}
Let
\[
V=\bigoplus_{\alpha\in I}V_\alpha
\]
be a finite direct sum over $K$, with packet projectors $E_\alpha$.  Let $\Lambda\subseteq V$ be an $\OO$ lattice.  Put
\[
\Lambda_\alpha=\Lambda\cap V_\alpha,
\qquad
\Lambda_{\mathrm{split}}=\bigoplus_\alpha\Lambda_\alpha,
\]
and define the gluing module
\[
G_{\Lambda,\Pcal}
=
\Lambda/\Lambda_{\mathrm{split}}.
\]
The separation weight of $E_\alpha$ is
\[
s_\Lambda(E_\alpha)
=
\min\{s\geq0:\pi^sE_\alpha\Lambda\subseteq\Lambda\}.
\]
The global separation weight is
\[
s_\Lambda(\Pcal)=\max_\alpha s_\Lambda(E_\alpha).
\]
For a finite torsion $\OO$ module $G$, write $\expn_\pi(G)$ for the least $s$ such that $\pi^sG=0$.
\end{definition}

The module $G_{\Lambda,\Pcal}$ is finite because $\Lambda_{\mathrm{split}}$ and $\Lambda$ are full lattices in the same vector space.

\begin{theorem}[Gluing module and exact separation weight]\label{thm:gluing_weight}
With the notation of \cref{def:lattice_packet}, the map
\[
\Phi:G_{\Lambda,\Pcal}
\longrightarrow
\bigoplus_{\alpha\in I}E_\alpha\Lambda/\Lambda_\alpha,
\qquad
[x]\longmapsto(E_\alpha x+\Lambda_\alpha)_\alpha
\]
is injective, and every coordinate projection of its image is surjective.  We also have
\[
\boxed{
\expn_\pi(G_{\Lambda,\Pcal})
=
s_\Lambda(\Pcal).
}
\]
If
\[
G_{\Lambda,\Pcal}
\cong
\bigoplus_{j=1}^r\OO/\pi^{a_j}\OO
\]
with $a_1\leq\cdots\leq a_r$, then
\[
s_\Lambda(\Pcal)=a_r.
\]
\end{theorem}

\begin{proof}
If $x\in\Lambda$ maps to zero, then $E_\alpha x\in\Lambda_\alpha$ for every $\alpha$.  Since
\[
x=\sum_\alpha E_\alpha x,
\]
we have $x\in\Lambda_{\mathrm{split}}$.  Thus $\Phi$ is injective.  The coordinate map onto $E_\alpha\Lambda/\Lambda_\alpha$ is surjective by the definition of $E_\alpha\Lambda$.

Let $s=s_\Lambda(\Pcal)$.  For every $x\in\Lambda$ and every $\alpha$,
\[
\pi^sE_\alpha x\in\Lambda_\alpha.
\]
Hence
\[
\pi^sx=\sum_\alpha\pi^sE_\alpha x\in\Lambda_{\mathrm{split}},
\]
so $\pi^sG_{\Lambda,\Pcal}=0$.  This proves
\[
\expn_\pi(G_{\Lambda,\Pcal})\leq s.
\]
If $d<s$, then for some $\alpha$ there is an $x\in\Lambda$ with
\[
\pi^dE_\alpha x\notin\Lambda_\alpha.
\]
The class $\pi^d[x]$ is therefore nonzero in the gluing module.  Thus $\pi^d$ does not annihilate the gluing module, and equality follows.  The final formula is the usual description of the exponent through Smith normal form.
\end{proof}

\begin{corollary}[Total gluing and worst precision]\label{cor:length_and_weight}
If
\[
G_{\Lambda,\Pcal}
\cong
\bigoplus_{j=1}^r\OO/\pi^{a_j}\OO,
\]
then
\[
\length_{\OO}(G_{\Lambda,\Pcal})
=
\sum_{j=1}^ra_j,
\qquad
s_\Lambda(\Pcal)=\max_ja_j.
\]
Thus the lattice index measures the total amount of gluing, while the separation weight measures the worst precision loss in one direction.
\end{corollary}

\begin{proof}
The length and exponent of a finite module are read directly from its Smith decomposition.  Apply \cref{thm:gluing_weight}.
\end{proof}

\begin{corollary}[Two complementary packets]\label{cor:binary_gluing}
Suppose
\[
V=V_1\oplus V_2
\]
with projector $E$ onto $V_1$ and complementary projector $I-E$.  Then
\[
\Lambda/(\Lambda_1\oplus\Lambda_2)
\cong
E\Lambda/\Lambda_1
\cong
(I-E)\Lambda/\Lambda_2.
\]
Consequently
\[
s_\Lambda(E)=s_\Lambda(I-E)
=
\expn_\pi\bigl(\Lambda/(\Lambda_1\oplus\Lambda_2)\bigr).
\]
\end{corollary}

\begin{proof}
Map $x\in\Lambda$ to $Ex+\Lambda_1$.  Its kernel consists of the $x$ for which $Ex\in\Lambda_1$.  In that case
\[
(I-E)x=x-Ex\in\Lambda\cap V_2=\Lambda_2,
\]
so the kernel is $\Lambda_1\oplus\Lambda_2$.  The map is surjective by definition.  The complementary map is handled in the same way.  The equality of weights follows from \cref{thm:gluing_weight}.
\end{proof}

\begin{corollary}[Integral splitting criterion]\label{cor:integral_split}
Every packet projector preserves $\Lambda$ if and only if
\[
\Lambda=\bigoplus_\alpha(\Lambda\cap V_\alpha).
\]
Equivalently,
\[
s_\Lambda(\Pcal)=0
\quad\Longleftrightarrow\quad
G_{\Lambda,\Pcal}=0.
\]
\end{corollary}

\begin{proof}
If every projector preserves $\Lambda$, then
\[
x=\sum_\alpha E_\alpha x
\]
expresses every $x\in\Lambda$ as a sum of vectors in $\Lambda_\alpha$.  The converse is immediate because each projector is the identity on one direct summand and zero on the others.  The final equivalence follows from \cref{thm:gluing_weight}.
\end{proof}

\subsection{Refinement and separation weight}

A finer rational packet system imposes more integral splitting conditions.

\begin{theorem}[Gluing exact sequence under refinement]\label{thm:gluing_refinement}
Let $\Qcal=(Q_\beta)$ refine $\Pcal=(P_\alpha)$ on $V$, and let $\Lambda$ be a lattice.  Define
\[
\Lambda_{\Pcal}
=
\bigoplus_\alpha(\Lambda\cap P_\alpha V),
\qquad
\Lambda_{\Qcal}
=
\bigoplus_\beta(\Lambda\cap Q_\beta V).
\]
Then
\[
\Lambda_{\Qcal}\subseteq\Lambda_{\Pcal}\subseteq\Lambda
\]
and there is a short exact sequence
\[
0
\longrightarrow
\Lambda_{\Pcal}/\Lambda_{\Qcal}
\longrightarrow
\Lambda/\Lambda_{\Qcal}
\longrightarrow
\Lambda/\Lambda_{\Pcal}
\longrightarrow0.
\]
In particular,
\[
\boxed{
s_\Lambda(\Pcal)\leq s_\Lambda(\Qcal).
}
\]
\end{theorem}

\begin{proof}
Every fine packet space lies in its parent packet space, so every element of $\Lambda_{\Qcal}$ belongs to $\Lambda_{\Pcal}$.  The displayed exact sequence is the standard quotient sequence for nested submodules.  The exponent of a quotient cannot exceed the exponent of the module from which it is obtained.  Apply \cref{thm:gluing_weight} to the two gluing modules.
\end{proof}

\begin{example}[Strict growth under refinement]\label{ex:strict_weight_growth}
Take the standard lattice $\mathbb Z_2^2$ and the involution from \cref{prop:sign_projectors}.  The one packet system $(I)$ has weight zero.  Refining it into the two rational sign packets produces weight one.  The coarse lattice is already split because there is only one summand.  The fine eigenspaces are glued together by one dyadic digit.
\end{example}

\begin{remark}[Strict growth can occur]\label{rem:strict_growth}
A coarse packet family may preserve the lattice even though a fine spectral split does not.  Refinement can reveal rational directions that are glued together integrally.  The weight records the extra depth needed to separate them.
\end{remark}

\subsection{Exact precision delay}

The separation weight has an operational meaning.  It is the exact number of extra input digits needed to recover a projected packet uniformly.

\begin{theorem}[Necessary and sufficient precision delay]\label{thm:precision_delay}
Let $E$ be a rational projector on $V$, let $\Lambda$ be a lattice, and let
\[
s=s_\Lambda(E).
\]
For every $m\geq0$, the rule
\[
[x]_{m+s}\longmapsto[Ex]_m
\]
from $\Lambda/\pi^{m+s}\Lambda$ to $V/\pi^m\Lambda$ is well defined.  If $d<s$, then for every $m\geq0$ there are $x,y\in\Lambda$ such that
\[
x\equiv y\pmod{\pi^{m+d}\Lambda}
\]
but
\[
Ex\not\equiv Ey\pmod{\pi^m\Lambda}.
\]
Thus $s$ extra digits are both sufficient and necessary for a uniform precision guarantee.
\end{theorem}

\begin{proof}
If
\[
x-y=\pi^{m+s}z,
\]
then
\[
E(x-y)=\pi^m(\pi^sEz)\in\pi^m\Lambda.
\]
This proves that the rule is well defined.  Now take $d<s$.  By the minimality of $s$, there is a $z\in\Lambda$ with
\[
\pi^dEz\notin\Lambda.
\]
Set
\[
x=\pi^{m+d}z,
\qquad
y=0.
\]
Then $x\equiv y$ modulo $\pi^{m+d}\Lambda$, while
\[
E(x-y)=\pi^m(\pi^dEz)\notin\pi^m\Lambda.
\]
\end{proof}

\begin{example}[The same rational split with two lattices]\label{ex:lattice_dependence}
Let
\[
V=Ke_1\oplus Ke_2
\]
and let $E$ project onto $Ke_1$.  On the split lattice
\[
\Lambda_0=\OO e_1\oplus\OO e_2,
\]
the weight is zero.  Fix $s\geq1$ and consider
\[
\Lambda_s
=
\OO e_1+
\OO(\pi^{-s}e_1+e_2).
\]
Then
\[
\Lambda_s\cap Ke_1=\OO e_1,
\qquad
\Lambda_s\cap Ke_2=\pi^s\OO e_2.
\]
The gluing module is isomorphic to $\OO/\pi^s\OO$, so
\[
s_{\Lambda_s}(E)=s.
\]
The rational decomposition has not changed.  Only the lattice has changed.
\end{example}

The example can be realized by an integral operator.  Define
\[
Ae_1=0,
\qquad
Ae_2=\pi^se_2.
\]
In the basis
\[
e_1,
\quad
\pi^{-s}e_1+e_2,
\]
the matrix of $A$ is
\[
\begin{pmatrix}
0&-1\\
0&\pi^s
\end{pmatrix}.
\]
Its eigenvalue factors have resultant $\pi^s$, and the exact projector weight on $\Lambda_s$ is also $s$.

\section{The dyadic sign split}\label{sec:dyadic}

The smallest example already shows why coefficient selectors and spectral averages must be kept distinct.

\begin{proposition}[Sign projectors over a lattice]\label{prop:sign_projectors}
Let
\[
A=
\begin{pmatrix}
1&1\\
0&-1
\end{pmatrix}
\]
act on $\mathbb Z_p^2$.  Over $\mathbb Q_p$, the sign projectors are
\[
E_+=\frac{I+A}{2}
=
\begin{pmatrix}
1&1/2\\
0&0
\end{pmatrix},
\]
\[
E_-=\frac{I-A}{2}
=
\begin{pmatrix}
0&-1/2\\
0&1
\end{pmatrix}.
\]
For every odd prime $p$, both projectors preserve $\mathbb Z_p^2$.  At $p=2$,
\[
s_2(E_+)=s_2(E_-)=1.
\]
The factors $x-1$ and $x+1$ have resultant $2$.
\end{proposition}

\begin{proof}
The displayed matrices are complementary idempotents.  The scalar $2$ is a unit in $\mathbb Z_p$ for odd $p$, so both matrices are integral there.  Over $\mathbb Z_2$, multiplying by $2$ clears the denominators, while neither projector preserves the standard lattice.  Thus both weights are one.  The resultant of the two linear factors is their root difference, which is $2$.
\end{proof}

Modulo $2$, the two eigenvalues become equal.  The rational sign packets remain distinct over $\mathbb Q_2$, but the standard lattice does not split into their intersections until one more dyadic digit is retained.

\begin{proposition}[Integral character averages away from the group order]\label{prop:character_integrality}
Let $\OO$ be a discrete valuation ring, let $N$ be a unit in $\OO$, and suppose $\OO$ contains a primitive $N$th root of unity $\omega$.  If an integral operator $U$ satisfies $U^N=I$, then
\[
E_j=\frac1N\sum_{k=0}^{N-1}\omega^{-jk}U^k
\]
preserves every $U$ stable $\OO$ lattice.  At primes dividing $N$, this formula alone gives no integral projector.
\end{proposition}

\begin{proof}
Every scalar coefficient in the displayed sum belongs to $\OO$ because $N$ is a unit and $\omega$ is integral.  Each power of $U$ preserves the lattice, so their $\OO$ linear combination does as well.  When $N$ is not a unit, division by $N$ is not an operation in $\OO$, so the same argument is unavailable.
\end{proof}

The final sentence does not assert that integral spectral packets never exist at a prime dividing $N$.  It says that they cannot be obtained from the unmodified character average without more structure.  Coefficient residue selectors remain exact because they never divide by the group order.

\section{Smith normal form and congruence lifting}\label{sec:smith}

The gluing module gives the exact common denominator of packet projectors.  Smith normal form gives the corresponding coordinate by coordinate solvability and branching rule.

\begin{standardfact}[Smith normal form over a discrete valuation ring]\label{fact:smith}
Let $M\in M_{m,n}(\OO)$.  There are invertible matrices
\[
U\in\operatorname{GL}_m(\OO),
\qquad
V\in\operatorname{GL}_n(\OO)
\]
such that
\[
UMV
=
\diag(\pi^{s_1},\ldots,\pi^{s_r},0,\ldots,0),
\]
where
\[
0\leq s_1\leq\cdots\leq s_r
\]
and $r$ is the rank over $K$.  The exponents are uniquely determined.  See \cite{Newman1972}.
\end{standardfact}

\begin{proposition}[Computing a packet weight from an inclusion matrix]\label{prop:weight_from_inclusion}
Choose an $\OO$ basis of $\Lambda$ and an $\OO$ basis of $\Lambda_{\mathrm{split}}$.  Let $B$ be the matrix of the inclusion
\[
\Lambda_{\mathrm{split}}\hookrightarrow\Lambda.
\]
If the Smith form of $B$ is
\[
\diag(\pi^{a_1},\ldots,\pi^{a_n}),
\qquad
0\leq a_1\leq\cdots\leq a_n,
\]
then
\[
G_{\Lambda,\Pcal}
\cong
\bigoplus_{j=1}^n\OO/\pi^{a_j}\OO,
\qquad
s_\Lambda(\Pcal)=a_n.
\]
\end{proposition}

\begin{proof}
The gluing module is the cokernel of the inclusion matrix $B$.  Smith normal form gives its invariant factors.  The weight formula follows from \cref{thm:gluing_weight}.
\end{proof}

\begin{theorem}[Smith weighted congruence solve]\label{thm:smith_congruence}
Let $M\in M_{m,n}(\OO)$ have Smith exponents $s_1,\ldots,s_r$.  Let the residue field have size $q$, fix $L\geq1$, and consider
\[
Mx\equiv b\pmod{\pi^L}.
\]
Put
\[
\widetilde b=Ub
\]
for a Smith matrix $U$ as in \cref{fact:smith}.  The congruence is solvable if and only if
\[
\widetilde b_i\equiv0
\pmod{\pi^{\min(s_i,L)}}
\qquad(1\leq i\leq r)
\]
and
\[
\widetilde b_i\equiv0
\pmod{\pi^L}
\qquad(r<i\leq m).
\]
When it is solvable, the number of solutions in $(\OO/\pi^L\OO)^n$ is
\[
\boxed{
q^{(n-r)L+\sum_{i=1}^r\min(s_i,L)}.
}
\]
\end{theorem}

\begin{proof}
Set $x=Vy$.  Multiplying the congruence by $U$ reduces it to independent equations
\[
\pi^{s_i}y_i\equiv\widetilde b_i\pmod{\pi^L}
\qquad(1\leq i\leq r)
\]
and zero equations in the remaining rows.  Multiplication by $\pi^{s_i}$ on $\OO/\pi^L\OO$ has image equal to the elements divisible by $\pi^{\min(s_i,L)}$ and kernel of size
\[
q^{\min(s_i,L)}.
\]
Each of the $n-r$ unused variables is free and contributes $q^L$ choices.  Multiplying the independent counts gives the formula.
\end{proof}

\begin{corollary}[Obstruction and branching data]\label{cor:smith_obstruction}
The transformed entries of $\widetilde b$ that fail the stated divisibility conditions are the complete finite congruence obstructions.  When all obstructions vanish, the Smith exponents determine the number of solution branches at depth $L$.
\end{corollary}

\begin{proof}
This is exactly the solvability and counting statement of \cref{thm:smith_congruence} written in invariant coordinates.
\end{proof}

\begin{example}[Three weighted directions]\label{ex:smith_three}
Take $p=2$, $L=6$, and the diagonal weights
\[
(0,1,3).
\]
The right side
\[
(1,2,8)
\]
is admissible.  The three coordinate kernels have sizes
\[
1,
\quad2,
\quad8,
\]
so there are $16$ solutions.  Replacing the final entry $8$ by $4$ violates the required divisibility by $2^3$ and makes the system unsolvable.
\end{example}

\section{Linear refinement defects}\label{sec:defects}

Exact coordinate refinement concerns two representations of the same object.  In applications one often chooses a new coarse operator rather than conjugating the fine operator exactly.  The difference is measured by an intertwining defect.

Let $X_f$ and $X_c$ be vector spaces or Banach spaces, let
\[
C:X_f\longrightarrow X_c,
\]
and let
\[
T_f\in\End(X_f),
\qquad
T_c\in\End(X_c).
\]
Define
\[
\Delta=CT_f-T_cC.
\]

\begin{theorem}[Power defect identity]\label{thm:power_defect}
For every integer $n\geq1$,
\[
CT_f^n-T_c^nC
=
\sum_{k=0}^{n-1}T_c^{n-1-k}\Delta T_f^k.
\]
\end{theorem}

\begin{proof}
The statement is clear for $n=1$.  For the induction step, write
\[
CT_f^{n+1}-T_c^{n+1}C
=
(CT_f^n-T_c^nC)T_f+T_c^n(CT_f-T_cC)
\]
and apply the induction hypothesis.
\end{proof}

\begin{corollary}[Power norm bound]\label{cor:power_norm}
For bounded operators between Banach spaces,
\[
\|CT_f^n-T_c^nC\|
\leq
\|\Delta\|
\sum_{k=0}^{n-1}
\|T_c\|^{n-1-k}\|T_f\|^k.
\]
\end{corollary}

\begin{proof}
Apply the triangle inequality and submultiplicativity to \cref{thm:power_defect}.
\end{proof}

\begin{corollary}[Polynomial defect]\label{cor:polynomial_defect}
For
\[
p(z)=\sum_{n=0}^Na_nz^n,
\]
we have
\[
Cp(T_f)-p(T_c)C
=
\sum_{n=1}^Na_n
\sum_{k=0}^{n-1}T_c^{n-1-k}\Delta T_f^k.
\]
\end{corollary}

\begin{proof}
The constant terms cancel because $CI=IC$.  Apply \cref{thm:power_defect} to every positive power and sum.
\end{proof}

\begin{theorem}[Resolvent defect identity]\label{thm:resolvent_defect}
If $z-T_f$ and $z-T_c$ are invertible, then
\[
C(z-T_f)^{-1}-(z-T_c)^{-1}C
=
(z-T_c)^{-1}\Delta(z-T_f)^{-1}.
\]
\end{theorem}

\begin{proof}
Multiply the claimed identity on the left by $z-T_c$ and on the right by $z-T_f$.  The left side becomes
\[
(z-T_c)C-C(z-T_f)
=
CT_f-T_cC
=
\Delta.
\]
Multiplying back by the two resolvents proves the formula.
\end{proof}

\begin{theorem}[Holomorphic functional calculus defect]\label{thm:holomorphic_defect}
Let $T_f$ and $T_c$ be bounded operators on complex Banach spaces.  Let $h$ be holomorphic on a neighborhood of both spectra, and let $\Gamma$ be a union of positively oriented contours surrounding both spectra inside that neighborhood.  Then
\[
Ch(T_f)-h(T_c)C
=
\frac{1}{2\pi i}
\int_\Gamma
h(z)(z-T_c)^{-1}\Delta(z-T_f)^{-1}\,dz.
\]
Consequently,
\begin{align*}
\|Ch(T_f)-h(T_c)C\|
&\leq
\frac{\operatorname{length}(\Gamma)}{2\pi}
\sup_{z\in\Gamma}|h(z)|\,\|\Delta\|\\
&\quad\times
\sup_{z\in\Gamma}\|(z-T_c)^{-1}\|
\sup_{z\in\Gamma}\|(z-T_f)^{-1}\|.
\end{align*}
\end{theorem}

\begin{proof}
The holomorphic functional calculus gives
\[
h(T)
=
\frac{1}{2\pi i}
\int_\Gamma h(z)(z-T)^{-1}\,dz.
\]
Subtract the two formulas after composing with $C$, then use \cref{thm:resolvent_defect} inside the integral.  The norm estimate follows from the standard contour bound.
\end{proof}

\begin{theorem}[Exponential defect identity]\label{thm:exponential_defect}
For bounded operators and every real or complex $t$,
\[
Ce^{tT_f}-e^{tT_c}C
=
\int_0^t e^{(t-s)T_c}\Delta e^{sT_f}\,ds,
\]
where the integral is taken along the straight segment from $0$ to $t$ when $t$ is complex.  For real $t\geq0$,
\[
\|Ce^{tT_f}-e^{tT_c}C\|
\leq
\|\Delta\|
\int_0^t
\|e^{(t-s)T_c}\|
\|e^{sT_f}\|\,ds.
\]
\end{theorem}

\begin{proof}
For real $t$, differentiate
\[
F(s)=e^{(t-s)T_c}Ce^{sT_f}.
\]
We obtain
\[
F'(s)=e^{(t-s)T_c}\Delta e^{sT_f}.
\]
Integration from $0$ to $t$ gives the identity.  The complex statement follows by the same calculation along the straight segment.  The norm bound follows from the integral formula.
\end{proof}

\subsection{Several refinement levels}

Let
\[
C_\ell:X_{\ell+1}\longrightarrow X_\ell
\]
for $0\leq\ell<L$, and let $T_\ell\in\End(X_\ell)$.  Put
\[
\Delta_\ell=C_\ell T_{\ell+1}-T_\ell C_\ell.
\]
For $a<b$, write
\[
C_{a,b}=C_aC_{a+1}\cdots C_{b-1},
\]
with $C_{a,a}=I$.

\begin{theorem}[Multilevel defect decomposition]\label{thm:multilevel_defect}
The total defect from level $L$ to level zero is
\[
C_{0,L}T_L-T_0C_{0,L}
=
\sum_{\ell=0}^{L-1}
C_{0,\ell}\Delta_\ell C_{\ell+1,L}.
\]
For bounded maps,
\[
\|C_{0,L}T_L-T_0C_{0,L}\|
\leq
\sum_{\ell=0}^{L-1}
\|C_{0,\ell}\|
\|\Delta_\ell\|
\|C_{\ell+1,L}\|.
\]
\end{theorem}

\begin{proof}
Insert the intermediate terms
\[
C_{0,\ell}T_\ell C_{\ell,L}
\]
for $0\leq\ell\leq L$.  Consecutive differences telescope.  The difference between the terms indexed by $\ell+1$ and $\ell$ is
\[
C_{0,\ell}(C_\ell T_{\ell+1}-T_\ell C_\ell)C_{\ell+1,L}.
\]
Summing gives the identity.  The norm estimate follows from the triangle inequality.
\end{proof}

\begin{example}[A visible one step defect]\label{ex:defect_example}
Let $X_f=\bbk^2$, let $X_c=\bbk$, and let
\[
C(x_1,x_2)=x_1.
\]
Take
\[
T_f=
\begin{pmatrix}
a&b\\
c&d
\end{pmatrix},
\qquad
T_c=(a).
\]
Then
\[
\Delta=CT_f-T_cC=(0,b).
\]
The coarse operator is exactly compatible with the fine operator precisely when the discarded coordinate does not feed into the retained coordinate.  The entry $c$ does not appear in the one step defect because it moves information in the opposite direction.  It can still affect later powers through the terms in \cref{thm:power_defect}.
\end{example}

\begin{corollary}[Exact compatibility]\label{cor:exact_compatibility}
If every one step defect $\Delta_\ell$ vanishes, then
\[
C_{0,L}p(T_L)=p(T_0)C_{0,L}
\]
for every polynomial $p$.  The analogous resolvent, holomorphic, and exponential identities also hold whenever the relevant expressions are defined.
\end{corollary}

\begin{proof}
The total one step defect vanishes by \cref{thm:multilevel_defect}.  Apply \cref{cor:polynomial_defect,thm:resolvent_defect,thm:holomorphic_defect,thm:exponential_defect}.
\end{proof}

\section{Scope and limits}\label{sec:limits}

The exact statements above have several deliberate boundaries.

First, refinement is a change of coordinates, not a truncation theorem.  If a numerical method discards fine packets, the resulting error must be estimated separately.

Second, the weighted inverse is linear.  A nonlinear map creates higher order interactions among packet coordinates.  Those interactions are not contained in the path formula of \cref{thm:weighted_inverse}.

Third, a resultant is only a denominator bound.  The exact weight depends on the lattice and is given by \cref{thm:gluing_weight}.  The two examples in \cref{ex:loose_resultant,ex:lattice_dependence} show that the same eigenvalue separation can give different integral packet weights.

Fourth, the dyadic sign example is a model rather than a classification of all phenomena at the prime two.  Character averages, trace splittings, quadratic refinements, and deformation projectors must each be checked on their own lattices.

Finally, the defect formulas do not prove that a chosen coarse operator is a good approximation.  They say exactly how a known one step defect propagates.  A useful approximation theorem must also bound that defect in the problem under study.


\bibliographystyle{aomplain}
\bibliography{PSC_Paper_II_Annals_Human_Edition}

\end{document}
